
%%%%%%%%%%%%%%%%%%%%%%%%%%%%%%%%%%%%%%%%%
% Developer CV
% LaTeX Class
% Version 2.0 (12/10/23)
%
% This class originates from:
% http://www.LaTeXTemplates.com
%
% Authors:
% Omar Roldan
% Based on a template by  Jan Vorisek (jan@vorisek.me)
% Based on a template by Jan Küster (info@jankuester.com)
% Modified for LaTeX Templates by Vel (vel@LaTeXTemplates.com)
%
% License:
% The MIT License (see included LICENSE file)
%
%%%%%%%%%%%%%%%%%%%%%%%%%%%%%%%%%%%%%%%%%

%----------------------------------------------------------------------------------------
%	PACKAGES AND OTHER DOCUMENT CONFIGURATIONS
%----------------------------------------------------------------------------------------

\documentclass[9pt]{developercv} % Default font size, values from 8-12pt are recommended
\usepackage{multicol}
\setlength{\columnsep}{0mm}
%----------------------------------------------------------------------------------------
\usepackage{lipsum}


\begin{document}

%----------------------------------------------------------------------------------------
%	TITLE AND CONTACT INFORMATION
%----------------------------------------------------------------------------------------

\begin{minipage}[t]{0.5\textwidth}
	\vspace{-\baselineskip} % Required for vertically aligning minipages

	{ \fontsize{16}{20} \textcolor{black}{\textbf{\MakeUppercase{Lorenzo Fumi}}}} % First name

	\vspace{6pt}

	{\Large Sviluppatore} % Career or current job title
\end{minipage}
\hfill
\begin{minipage}[t]{0.2\textwidth} % 20% of the page width for the first row of icons
	\vspace{-\baselineskip} % Required for vertically aligning minipages

	% The first parameter is the FontAwesome icon name, the second is the box size and the third is the text
	%\icon{Globe}{11}{\href{http://www.google.com}{portafolio.com}}\\
    \icon{Phone}{11}{+39 328 338 8037}\\
    \icon{MapMarker}{11}{Piacenza, Italia}\\

\end{minipage}
\begin{minipage}[t]{0.27\textwidth} % 27% of the page width for the second row of icons
	\vspace{-\baselineskip} % Required for vertically aligning minipages

	\icon{Envelope}{11}{\href{mailto:lorenzofumi@gmail.com}{lorenzofumi@gmail.com}}\\
    \icon{Github}{11}{\href{https://github.com/DeeJack}{github.com/DeeJack}}\\
    \icon{LinkedinSquare}{11}{\href{https://www.linkedin.com/in/lorenzo-fumi-017471222/}{Lorenzo Fumi}}\\

\end{minipage}


%----------------------------------------------------------------------------------------
%	INTRODUCTION, SKILLS AND TECHNOLOGIES
%----------------------------------------------------------------------------------------

\begin{minipage}[t]{0.46\textwidth}
    \cvsect{Riassunto}
	\vspace{-6pt}

    Neolaureato magistrale in Informatica con tesi su sistemi di agenti LLM. Esperienza nello sviluppo software, Data Science, e AI. Particolare interesse in sistemi di agenti AI.\\
\end{minipage}
\hfill % Whitespace between
\begin{minipage}[t]{0.465\textwidth}
    \cvsect{Skills}
    \vspace{-6pt}

    \begin{minipage}[t]{0.2\textwidth}
        \textbf{Linguaggi:}
    \end{minipage}
    \hfill
    \begin{minipage}[t]{0.73\textwidth}
      Python, Java, SQL, JavaScript, HTML, CSS, C\#.
    \end{minipage}
    \vspace{4mm}

    \begin{minipage}[t]{0.2\textwidth}
        \textbf{Librerie:}
    \end{minipage}
    \hfill
    \begin{minipage}[t]{0.73\textwidth}
      LangChain, LangGraph, Pandas, NumPy, scikit-learn.
    \end{minipage}
    \vspace{4mm}

    \begin{minipage}[t]{0.2\textwidth}
        \textbf{Database:}
    \end{minipage}
    \hfill
    \begin{minipage}[t]{0.73\textwidth}
      MS SQL Server, MongoDB, MySQL.
    \end{minipage}
    \vspace{4mm}

    \begin{minipage}[t]{0.2\textwidth}
        \textbf{Tecnologie:}
    \end{minipage}
    \hfill
    \begin{minipage}[t]{0.73\textwidth}
      GIT, Docker, Node.js, ASP.NET, Angular, Vue.js.
    \end{minipage}

\end{minipage}

%----------------------------------------------------------------------------------------
%	Projects
%----------------------------------------------------------------------------------------
\cvsect{Progetti}
\begin{entrylist}
    \entry
		{LLM Agents}
		{ReflectOR: Agentic System for Post-Operative Surgical Debriefing}
		{}
		{Sistema di agenti LLM per assistere i chirurghi in debriefing post-operatori, e creare documentazioni per utilizzo futuro. Il sistema trascrive e diarizza operazioni chirurgiche, le analizza con agenti LLM, e supporta un debriefing strutturato. Progetto tesi magistrale.

        Tecnologie: Python, LangChain, LangGraph}
	\entry
		{ML}
		{iRanking}
		{}
		{Sistema che usa machine learning per misurare l'impatto della temperatura sul tempo di completamento in un simulatore (iRacing). Sviluppo di modelli di Machine Learning e analisi dei dati.

        Tecnologie: Python, scikit-learn, Pandas, NumPy}
    \entry
		{Android}
		{Roommates}
		{}
		{App Android sviluppata per un corso universitario per organizzare compiti e spese tra coinquilini.

        Tecnologie: Java}
    \entry
		{Personal}
		{Diary}
		{}
		{Applicazione per prendere note con Surface Pen, sviluppata in PyQt6.

        Tecnologie: Python, PyQt6}
\end{entrylist}

%----------------------------------------------------------------------------------------
%	EDUCATION
%----------------------------------------------------------------------------------------
\vspace{-10 pt}
\cvsect{Educazione}
\begin{entrylist}
    \entry
		{9/2020 - 9/2023}
		{Laurea triennale in Informatica}
		{Università di Trento (1° anno presso Università di Parma)}
		{Voto: 110 e Lode. Tesi: ``iRanking: Raffinamento del Sistema di Classifica di un Simulatore Online'', uno studio dell'impatto della temperatura sul tempo di completamento dei giri su un simulatore di macchine.}
	\entry
		{9/2023 - 10/2025}
		{Laurea magistrale in Informatica}
		{Università di Trento}
		{Voto: 110 e Lode. Tesi: ``ReflectOR: Agentic System for Post-Operative Surgical Debriefing'', sistema di agenti LLM applicato sul dominio del debriefing post-operatorio}
\end{entrylist}

%----------------------------------------------------------------------------------------
%	EXPERIENCE
%----------------------------------------------------------------------------------------
\vspace{-10 pt}
\cvsect{Esperienze}
\begin{entrylist}
	\entry
		{7/2018 -- 8/2018}
		{Stage estivo (2 mesi)}
		{H\&S}
		{\vspace{-10pt}
        \begin{itemize}[noitemsep,topsep=0pt,parsep=0pt,partopsep=0pt, leftmargin=-1pt]
            \item Sviluppo con C# in ASP.NET e JavaScript per front-end e back-end.
            \item Correzione bug e aggiunta di nuove funzionalità.
        \end{itemize}
        \texttt{C\#} \slashsep \texttt{ASP.NET} \slashsep \texttt{JavaScript}}
	\entry
        {5/2019 -- 5/2019}
		{Alternanza Scuola-Lavoro (1 mese)}
		{Zucchetti}
		{\vspace{-10pt}
        \begin{itemize}[noitemsep,topsep=0pt,parsep=0pt,partopsep=0pt, leftmargin=-1pt]
            \item Utilizzo del programma Kettle per migrare i dati da un database IBM DB2 a MS SQL Server.
            \item Sviluppo di query SQL e script per la gestione dei formati.
        \end{itemize}
        \texttt{SQL} \slashsep \texttt{Kettle}}
	\entry
		{10/2019 -- 10/2019}
		{Alternanza Scuola-Lavoro (1 mese)}
		{Techsol}
		{\vspace{-10pt}
        \begin{itemize}[noitemsep,topsep=0pt,parsep=0pt,partopsep=0pt, leftmargin=-1pt]
            \item Sviluppo di un'applicazione Angular per la gestione delle ferie dei dipendenti.
            \item Implementazione da zero di front-end e back-end.
        \end{itemize}
        \texttt{Angular} \slashsep \texttt{HTML} \slashsep \texttt{CSS}}
\end{entrylist}

%----------------------------------------------------------------------------------------
%	LANGUAGES
%----------------------------------------------------------------------------------------
\vspace{-10 pt}
	\cvsect{Lingue}
    \vspace{-6pt}

    \hspace{26mm} \textbf{ Italiano} - madrelingua, \textbf{English} - B2

%----------------------------------------------------------------------------------------

\end{document}
